\subsection{Sixteen calibrated formula encodings from one native calculus}
\label{sec:constructor-orientation}
There is a separate way to trade source operations without weakening the
logic: reverse the physical pair order independently for each formula
constructor.  The ordinary stack and proof-list pairing can remain unchanged.
The general all-code constructor-transport method and the separate
domination argument are inherited from the pinned zero-spare-transfer
report's prefix-recoding section~\cite{recoding}.  Here they are specialized to the
sixteen operand orientations and independently checked; the general
method is not claimed as new.

Write $J_r(a,b)=4P(a,b)+r$ for tags $r=0,1,2,3$, respectively implication,
membership, universal quantification and equality.  For an arbitrary bit
vector $e\in\{0,1\}^4$, put
\[
 J^e_r(a,b)=
 \begin{cases}
 4P(a,b)+r,&e_r=0,\\
 4P(b,a)+r,&e_r=1.
 \end{cases}
\]
All constructors, including constructor-valued numerical labels inside
other constructors, use this same chosen interpretation.  Truth and falsity
remain the literal numbers zero and one.

\paragraph{All-natural isomorphism.}
Define $F_e(0)=0$.  For $n>0$, write uniquely $n=4P(a,b)+r$ and set
\[
 F_e(n)=J^e_r(F_e(a),F_e(b)).
\]
Both children are strictly below $n$, including the cases $n=1,2,3$,
whose two children are zero.  This is primitive-recursive
course-of-values recursion.  For every natural $a,b$ and tag $r$,
\[
 F_e(J_r(a,b))=J^e_r(F_e(a),F_e(b)).
\]
The only zero-input case is $J_0(0,0)=0$, where the equation still holds.
The map fixes one because $1=J_1(0,0)$.  Strong induction on $n$ now proves
$F_e(F_e(n))=n$: each tag's swap occurs twice, and both child subcodes
satisfy the induction hypothesis.  Thus $F_e$ is an involution on
\emph{all} natural codes, not only conventionally well-formed formulas.

The unrestricted native calculi are exactly isomorphic under $F_e$.
Every schema and theorem-only generalization step transports by the
constructor identity; numerical MP transports by exact equality.
Numerical labels transform along with formula codes.  The full bundle's
gate $C\ne A,B$ is preserved by injectivity, and its literal falsity one
is fixed.  Hence
\[
 \mathcal C\vdash n\quad\Longleftrightarrow\quad
 \mathcal C_e\vdash F_e(n).
\]

\paragraph{The dominance bound requires a separate argument.}
The map does not preserve numerical order: reversing implication exchanges
the codes $I(1,0)=4$ and $I(0,1)=8$.  It would therefore be incorrect to
transport $U\le A$ through $F_e$ without further work.

Instead interpret the accepted 3,728-node dominated-row certificate using
the new implication constructor.  Its1,260 eighth-row leaves have fourth
parameter either zero (224 leaves), the entire first parameter (164), or
an implication subterm of the first parameter (872).  Thus every bound
uses only
\[
 I(0,0)=0,\qquad I(a,b)\ge a,\qquad I(a,b)\ge b.
\]
Both physical orientations satisfy these identities and inequalities.
The exact MP equations and zero normalization in the certificate also
remain valid.  Consequently the oriented dominated calculus
$\mathcal D_e$ recovers its unrestricted eighth row and has exactly the
same theorem codes as $\mathcal C_e$.  In particular,
\[
 \mathcal D_e\vdash1
 \ \Longleftrightarrow\ 
 \mathcal C_e\vdash1
 \ \Longleftrightarrow\ 
 \mathcal C\vdash1
 \ \Longleftrightarrow\ 
 \operatorname{Incon}(\mathrm{ZFC}),
\]
where the final comparison is the inherited native-calculus calibration.

This supplies sixteen safe encoding choices for source optimization.
An actual edit must establish the selected constructor consistently at
every formula call while preserving the stack and scheduler interfaces.
An indiscriminate reversal of a shared pairing helper need not do that.
Physical allocation, primitive simulation, and blank initialization
remain separate obligations.
